\begin{ex}
Trong hai phân số $\frac{3}{7}$ và $\frac{5}{7}$, phân số nào lớn hơn?
\shortans{5/7}
\loigiai{
Cùng mẫu số 7, $5 > 3$ nên $\frac{5}{7} > \frac{3}{7}$.
}
\end{ex}

\begin{ex}
Trong hai phân số $\frac{8}{11}$ và $\frac{6}{11}$, phân số nào bé hơn?
\shortans{6/11}
\loigiai{
$\frac{6}{11} < \frac{8}{11}$.
}
\end{ex}

\begin{ex}
Trong hai phân số có cùng tử số $\frac{3}{5}$ và $\frac{3}{7}$, phân số nào lớn hơn?
\shortans{3/5}
\loigiai{
Cùng tử số, mẫu số bé hơn thì phân số lớn hơn: $\frac{3}{5} > \frac{3}{7}$.
}
\end{ex}

\begin{ex}
Trong hai phân số $\frac{4}{9}$ và $\frac{4}{7}$, phân số nào bé hơn?
\shortans{4/9}
\loigiai{
Mẫu số lớn hơn thì bé hơn: $\frac{4}{9} < \frac{4}{7}$.
}
\end{ex}

\begin{ex}
So sánh $\frac{2}{3}$ và $\frac{3}{4}$: phân số nào lớn hơn?
\shortans{3/4}
\loigiai{
$\frac{8}{12} < \frac{9}{12}$ nên $\frac{3}{4}$ lớn hơn.
}
\end{ex}

\begin{ex}
So sánh $\frac{5}{6}$ và $\frac{4}{5}$: phân số nào lớn hơn?
\shortans{5/6}
\loigiai{
$\frac{25}{30} > \frac{24}{30}$ nên $\frac{5}{6}$ lớn hơn.
}
\end{ex}

\begin{ex}
So sánh $\frac{7}{8}$ và 1: phân số $\frac{7}{8}$ bé hơn hay lớn hơn 1? (Bé hơn ghi 0, lớn hơn ghi 1).
\shortans{0}
\loigiai{
Tử < Mẫu nên bé hơn 1 (0).
}
\end{ex}

\begin{ex}
So sánh $\frac{9}{7}$ và 1: phân số $\frac{9}{7}$ lớn hơn hay bé hơn 1? (Lớn hơn ghi 1, bé hơn ghi 0).
\shortans{1}
\loigiai{
Tử > Mẫu nên lớn hơn 1 (1).
}
\end{ex}

\begin{ex}
So sánh $\frac{5}{5}$ và 1: hai giá trị này có bằng nhau không? (Bằng nhau ghi 1, khác nhau ghi 0).
\shortans{1}
\loigiai{
$\frac{5}{5} = 1$ (1).
}
\end{ex}

\begin{ex}
Trong các phân số $\frac{1}{2}, \frac{2}{3}, \frac{3}{4}, \frac{4}{5}$, phân số nào lớn nhất?
\shortans{4/5}
\loigiai{
Quy đồng mẫu 60: 30/60, 40/60, 45/60, 48/60. Lớn nhất là $\frac{4}{5}$.
}
\end{ex}

\begin{ex}
Trong các phân số trên, phân số nào bé nhất?
\shortans{1/2}
\loigiai{
$\frac{1}{2}$ là bé nhất.
}
\end{ex}

\begin{ex}
Sắp xếp theo thứ tự tăng dần: trong $\frac{2}{5}, \frac{4}{5}, \frac{1}{5}, \frac{3}{5}$, phân số lớn thứ nhì là phân số nào?
\shortans{3/5}
\loigiai{
Thứ tự: 1/5, 2/5, 3/5, 4/5. Lớn thứ nhì là $\frac{3}{5}$.
}
\end{ex}

\begin{ex}
So sánh $\frac{1}{4}$ và $\frac{2}{8}$: hai phân số này có bằng nhau không? (Bằng nhau ghi 1, khác nhau ghi 0).
\shortans{1}
\loigiai{
$\frac{2}{8} = \frac{1}{4}$ (bằng nhau).
}
\end{ex}

\begin{ex}
So sánh $\frac{3}{5}$ và $\frac{6}{10}$: hai phân số này có bằng nhau không? (Bằng nhau ghi 1, khác nhau ghi 0).
\shortans{1}
\loigiai{
Bằng nhau (1).
}
\end{ex}

\begin{ex}
Hiệu giữa 1 và $\frac{3}{4}$ là bao nhiêu?
\shortans{1/4}
\loigiai{
$1 - \frac{3}{4} = \frac{1}{4}$.
}
\end{ex}

\begin{ex}
Hiệu giữa 1 và $\frac{5}{6}$ là bao nhiêu?
\shortans{1/6}
\loigiai{
$1 - \frac{5}{6} = \frac{1}{6}$.
}
\end{ex}

\begin{ex}
Phân số nào gần 1 hơn: $\frac{7}{8}$ hay $\frac{8}{9}$?
\shortans{8/9}
\loigiai{
$1 - \frac{8}{9} = \frac{1}{9} < \frac{1}{8}$ nên $\frac{8}{9}$ gần 1 hơn.
}
\end{ex}

\begin{ex}
Trong hai phân số $\frac{11}{12}$ và $\frac{12}{13}$, phân số nào lớn hơn?
\shortans{12/13}
\loigiai{
So sánh phần bù: $1/13 < 1/12$ nên $\frac{12}{13} > \frac{11}{12}$.
}
\end{ex}

\begin{ex}
So sánh $\frac{15}{14}$ và $\frac{13}{14}$: phân số nào lớn hơn?
\shortans{15/14}
\loigiai{
$\frac{15}{14} > \frac{13}{14}$.
}
\end{ex}

\begin{ex}
Phân số nào bé hơn 1 trong hai phân số: $\frac{5}{6}$ và $\frac{6}{5}$?
\shortans{5/6}
\loigiai{
$\frac{5}{6} < 1$.
}
\end{ex}

\begin{ex}
Phân số nào lớn hơn 1 trong hai phân số: $\frac{8}{7}$ và $\frac{7}{8}$?
\shortans{8/7}
\loigiai{
$\frac{8}{7} > 1$.
}
\end{ex}

\begin{ex}
So sánh $\frac{3}{8}$ và $\frac{1}{2}$: phân số nào lớn hơn?
\shortans{1/2}
\loigiai{
$\frac{1}{2} = \frac{4}{8} > \frac{3}{8}$.
}
\end{ex}

\begin{ex}
So sánh $\frac{5}{8}$ và $\frac{1}{2}$: phân số nào lớn hơn?
\shortans{5/8}
\loigiai{
$\frac{5}{8} > \frac{4}{8} = \frac{1}{2}$.
}
\end{ex}

\begin{ex}
Trong các phân số $\frac{2}{7}, \frac{5}{7}, \frac{3}{7}, \frac{6}{7}$, số lớn nhất là số nào?
\shortans{6/7}
\loigiai{
$\frac{6}{7}$ là lớn nhất.
}
\end{ex}

\begin{ex}
Số bé nhất trong các phân số trên là số nào?
\shortans{2/7}
\loigiai{
$\frac{2}{7}$ là bé nhất.
}
\end{ex}
